\documentclass{article}
\usepackage{amsmath,amssymb,amsthm}
\usepackage{geometry}
\geometry{margin=1in}
\newtheorem{theorem}{Theorem}
\newtheorem{lemma}{Lemma}
\newcommand{\prox}{\operatorname{prox}}
\newcommand{\ip}[2]{\langle #1,#2\rangle}
\newcommand{\norm}[1]{\left\|#1\right\|}
\begin{document}

% -------------------------------------------------
\section*{Primal–Dual Hybrid Gradient (PDHG)}
% -------------------------------------------------

\section*{Mathematical Formulation}
Consider the convex–concave saddle–point problem  

\[
\min_{x\in\mathbb{R}^{n}} \; f(x)+g(Kx)
\qquad\Longleftrightarrow\qquad
\min_{x}\max_{y}\; \mathcal{L}(x,y):=f(x)+\langle Kx,y\rangle-g^{*}(y),
\]

where  

\begin{itemize}
  \item $f:\mathbb{R}^{n}\!\to\!\mathbb{R}\cup\{+\infty\}$ is proper, convex and (possibly) non‑smooth,
  \item $g:\mathbb{R}^{m}\!\to\!\mathbb{R}\cup\{+\infty\}$ is proper, convex and (possibly) non‑smooth,
  \item $K\in\mathbb{R}^{m\times n}$ is a linear operator,
  \item $g^{*}$ denotes the Fenchel conjugate of $g$.
\end{itemize}

The PDHG iterates $(x_{k},y_{k})$ are defined by  

\[
\boxed{
\begin{aligned}
y_{k+1} &= \prox_{\sigma g^{*}}\!\bigl(y_{k}+\sigma K\bar{x}_{k}\bigr),\\[4pt]
x_{k+1} &= \prox_{\tau f}\!\bigl(x_{k}-\tau K^{\!\top}y_{k+1}\bigr),\\[4pt]
\bar{x}_{k+1} &= x_{k+1}+\theta\,(x_{k+1}-x_{k}),
\end{aligned}}
\tag{PDHG}
\]

with step sizes $\tau,\sigma>0$, extrapolation parameter $\theta\in[0,1]$, and  

\[
\bar{x}_{0}=x_{0},\qquad
\text{where } \prox_{\lambda h}(v):=\arg\min_{u}\Big\{\frac{1}{2\lambda}\|u-v\|^{2}+h(u)\Big\}.
\]

The parameters must satisfy  

\[
\tau\sigma\|K\|^{2}<1,
\tag{A1}
\]

where $\|K\|$ is the spectral norm of $K$.

\section*{Pseudocode}
\begin{verbatim}
Input:   f, g, linear operator K
         step sizes τ>0, σ>0, extrapolation θ∈[0,1]
         initial points x0, y0
Initialize: x = x0, y = y0, x_bar = x0
for k = 0,1,2,...,T-1 do
    y = prox_{σ g*}( y + σ * K * x_bar )
    x_new = prox_{τ f}( x - τ * K^T * y )
    x_bar = x_new + θ * (x_new - x)
    x = x_new
end for
Output: (x, y)
\end{verbatim}

\section*{Dry Run Example}
We illustrate three iterations on the scalar problem  

\[
\min_{w\in\mathbb{R}} \;(w-3)^{2},
\]

which can be written as the saddle problem with  

\[
f(w)=(w-3)^{2},\qquad g\equiv0,\qquad K=1.
\]

Thus $g^{*}=0$ and $\prox_{\sigma g^{*}}(u)=u$.  
Choose $\tau=\sigma=0.1$, $\theta=1$, and initialise $w_{0}=0$, $y_{0}=0$.
The proximal operator of $f$ is computed analytically:

\[
\prox_{\tau f}(v)=\arg\min_{w}\Big\{\frac{1}{2\tau}(w-v)^{2}+(w-3)^{2}\Big\}
               =\frac{v/\tau+6}{1/\tau+2}.
\tag{1}
\]

\medskip
\noindent\textbf{Iteration $k=0$}
\[
\begin{aligned}
\bar{w}_{0}&=w_{0}=0,\\
y_{1}&=y_{0}+\sigma\bar{w}_{0}=0+0.1\cdot0=0,\\
w_{1}&=\prox_{\tau f}\!\bigl(w_{0}-\tau y_{1}\bigr)
      =\prox_{0.1 f}(0-0)=\frac{0/0.1+6}{1/0.1+2}
      =\frac{0+6}{10+2}=0.5,\\
\bar{w}_{1}&=w_{1}+\theta(w_{1}-w_{0})=0.5+1\cdot(0.5-0)=1.0.
\end{aligned}
\]
Objective value $f(w_{1})=(0.5-3)^{2}=6.25$.

\medskip
\noindent\textbf{Iteration $k=1$}
\[
\begin{aligned}
y_{2}&=y_{1}+\sigma\bar{w}_{1}=0+0.1\cdot1.0=0.1,\\
w_{2}&=\prox_{\tau f}\!\bigl(w_{1}-\tau y_{2}\bigr)
      =\prox_{0.1 f}(0.5-0.01)=\prox_{0.1 f}(0.49)\\
     &=\frac{0.49/0.1+6}{10+2}
      =\frac{4.9+6}{12}=0.9083\;( \text{rounded}),\\
\bar{w}_{2}&=w_{2}+ (w_{2}-w_{1})
            =0.9083+(0.9083-0.5)=1.3166,\\
f(w_{2})&\approx(0.9083-3)^{2}=4.376.
\end{aligned}
\]

\medskip
\noindent\textbf{Iteration $k=2$}
\[
\begin{aligned}
y_{3}&=y_{2}+\sigma\bar{w}_{2}=0.1+0.1\cdot1.3166=0.23166,\\
w_{3}&=\prox_{\tau f}\!\bigl(w_{2}-\tau y_{3}\bigr)
      =\prox_{0.1 f}(0.9083-0.023166)=\prox_{0.1 f}(0.88513)\\
     &=\frac{0.88513/0.1+6}{12}
      =\frac{8.8513+6}{12}=1.2376,\\
\bar{w}_{3}&=w_{3}+ (w_{3}-w_{2})
            =1.2376+(1.2376-0.9083)=1.5669,\\
f(w_{3})&\approx(1.2376-3)^{2}=3.051.
\end{aligned}
\]

The iterates approach the optimum $w^{*}=3$, and the objective values decrease.

% -------------------------------------------------
\section*{Assumptions}
\begin{enumerate}
  \item[(A1)] \textbf{Step‑size condition:} $\tau\sigma\|K\|^{2}<1$.
  \item[(A2)] \textbf{Convexity:} $f$ and $g$ are proper, closed, convex.
  \item[(A3)] \textbf{Existence of a saddle point:} there exists $(x^{*},y^{*})$ such that  
        \[
        \mathcal{L}(x^{*},y)\le \mathcal{L}(x^{*},y^{*})\le \mathcal{L}(x,y^{*}),
        \quad\forall x,y.
        \]
  \item[(A4)] \textbf{Boundedness of iterates (optional):} $\{(x_{k},y_{k})\}$ is contained in a bounded set.
\end{enumerate}

% -------------------------------------------------
\section*{Main Convergence Theorem}
\begin{theorem}[Ergodic $O(1/T)$ convergence]\label{thm:main}
Under (A1)–(A3) the averaged iterates  

\[
\bar{x}_{T}:=\frac{1}{T}\sum_{k=1}^{T}x_{k},\qquad 
\bar{y}_{T}:=\frac{1}{T}\sum_{k=1}^{T}y_{k}
\]

satisfy the primal–dual gap bound  

\[
\boxed{
\mathcal{G}_{T}:=
\sup_{y}\bigl\{\mathcal{L}(\bar{x}_{T},y)-\mathcal{L}(x^{*},\bar{y}_{T})\bigr\}
\le
\frac{C}{T},
}
\tag{2}
\]

where the constant  

\[
C:=\frac{1}{2\tau}\norm{x_{0}-x^{*}}^{2}
   +\frac{1}{2\sigma}\norm{y_{0}-y^{*}}^{2}.
\tag{3}
\]

Consequently $\mathcal{L}(\bar{x}_{T},\bar{y}_{T})\to \mathcal{L}(x^{*},y^{*})$ at rate $O(1/T)$.
\end{theorem}

% -------------------------------------------------
\section*{Theoretical Properties}
\begin{itemize}
  \item \textbf{Stability:} The condition $\tau\sigma\|K\|^{2}<1$ guarantees that the coupled proximal updates are non‑expansive, yielding a Lyapunov function that monotonically decreases.
  \item \textbf{Hyper‑parameter sensitivity:} The bound $C$ in \eqref{eq:2} depends linearly on $1/\tau$ and $1/\sigma$. Larger step sizes (still satisfying (A1)) reduce $C$ but may violate the contraction condition.
  \item \textbf{Comparison to baselines:}
    \begin{itemize}
      \item Compared with vanilla gradient descent on smooth $f$, PDHG attains the same $O(1/T)$ rate without requiring differentiability of $f$ or $g$.
      \item Against ADMM, PDHG enjoys a simpler per‑iteration cost (only two proximal evaluations) and the same ergodic rate under comparable assumptions.
    \end{itemize}
\end{itemize}

% -------------------------------------------------
\section*{Discussion}
The theorem shows that PDHG inherits the optimal $O(1/T)$ rate for convex–concave saddle‑point problems while handling non‑smooth terms via proximal maps. The ergodic averaging is essential for the rate; without averaging one can still obtain $o(1)$ convergence of the primal residual under stronger assumptions (e.g., strong convexity). The step‑size condition (A1) is tight: if $\tau\sigma\|K\|^{2}=1$ the Lyapunov argument collapses, and divergence may occur.

% -------------------------------------------------
\section*{Preliminaries (Definitions and Lemmas)}
\begin{lemma}[Three‑point identity for proximal operators]\label{lem:three}
For any convex $h$ and $\lambda>0$, let $u=\prox_{\lambda h}(v)$. Then for all $z$,
\[
\langle u-v,\,z-u\rangle \ge \lambda\bigl( h(u)-h(z)\bigr).
\tag{4}
\]
\end{lemma}
\begin{proof}
The optimality condition of the proximal subproblem gives  
$0\in \frac{1}{\lambda}(u-v)+\partial h(u)$, i.e.
$\frac{v-u}{\lambda}\in\partial h(u)$. By monotonicity of $\partial h$,
$\langle \frac{v-u}{\lambda},z-u\rangle\ge h(z)-h(u)$. Multiplying by $\lambda$ yields \eqref{eq:4}.  \qedhere
\end{proof}

\begin{lemma}[Fenchel–Young inequality]\label{lem:fy}
For any $a\in\mathbb{R}^{m},b\in\mathbb{R}^{m}$,
\[
\langle a,b\rangle\le g(a)+g^{*}(b).
\tag{5}
\]
\end{lemma}
\begin{proof}
Immediate from the definition of the conjugate $g^{*}(b)=\sup_{a}\{\langle a,b\rangle-g(a)\}$. \qedhere
\end{proof}

% -------------------------------------------------
\section*{Main Proof (Step‑by‑Step Derivation)}
We prove Theorem~\ref{thm:main}.  
All steps are numbered for easy reference.

\subsection*{Notation}
Denote the saddle point by $(x^{*},y^{*})$.  
Define the Lyapunov function  

\[
\Phi_{k}:=
\frac{1}{2\tau}\norm{x_{k}-x^{*}}^{2}
+\frac{1}{2\sigma}\norm{y_{k}-y^{*}}^{2}.
\tag{6}
\]

\subsection*{Derivation}
\begin{enumerate}
\item[Step 1.] Apply Lemma~\ref{lem:three} to the $y$‑update:
\[
\langle y_{k+1}- (y_{k}+\sigma K\bar{x}_{k}),\;
      y^{*}-y_{k+1}\rangle
      \ge \sigma\bigl(g^{*}(y_{k+1})-g^{*}(y^{*})\bigr).
\tag{7}
\]
Rearranging gives
\[
\langle y_{k+1}-y_{k},\,y^{*}-y_{k+1}\rangle
      +\sigma\langle K\bar{x}_{k},\,y^{*}-y_{k+1}\rangle
      \ge \sigma\bigl(g^{*}(y_{k+1})-g^{*}(y^{*})\bigr).
\tag{8}
\]

\item[Step 2.] Apply Lemma~\ref{lem:three} to the $x$‑update:
\[
\langle x_{k+1}- (x_{k}-\tau K^{\!\top}y_{k+1}),\;
      x^{*}-x_{k+1}\rangle
      \ge \tau\bigl(f(x_{k+1})-f(x^{*})\bigr).
\tag{9}
\]
After expanding,
\[
\langle x_{k+1}-x_{k},\,x^{*}-x_{k+1}\rangle
      -\tau\langle K^{\!\top}y_{k+1},\,x^{*}-x_{k+1}\rangle
      \ge \tau\bigl(f(x_{k+1})-f(x^{*})\bigr).
\tag{10}
\]

\item[Step 3.] Add \eqref{eq:8} and \eqref{eq:10} and use the identity
$\langle a,b\rangle =\frac12(\|a\|^{2}+\|b\|^{2}-\|a-b\|^{2})$ to obtain
\[
\begin{aligned}
&\frac{1}{2\tau}\bigl(\|x_{k}-x^{*}\|^{2}-\|x_{k+1}-x^{*}\|^{2}
                     -\|x_{k+1}-x_{k}\|^{2}\bigr)\\
&\quad+\frac{1}{2\sigma}\bigl(\|y_{k}-y^{*}\|^{2}-\|y_{k+1}-y^{*}\|^{2}
                     -\|y_{k+1}-y_{k}\|^{2}\bigr)\\
&\quad+\langle K\bar{x}_{k},y^{*}-y_{k+1}\rangle
      -\langle K^{\!\top}y_{k+1},x^{*}-x_{k+1}\rangle  \\
&\ge f(x_{k+1})-f(x^{*})+g^{*}(y_{k+1})-g^{*}(y^{*}).
\end{aligned}
\tag{11}
\]

\item[Step 4.] Observe that
\[
\langle K\bar{x}_{k},y^{*}-y_{k+1}\rangle
-\langle K^{\!\top}y_{k+1},x^{*}-x_{k+1}\rangle
=
\langle K\bar{x}_{k},y^{*}\rangle
-\langle Kx^{*},y_{k+1}\rangle
-\langle K(x_{k+1}-\bar{x}_{k}),y_{k+1}\rangle.
\tag{12}
\]

\item[Step 5.] Use the definition of the saddle function
$\mathcal{L}(x,y)=f(x)+\langle Kx,y\rangle-g^{*}(y)$ to rewrite the right‑hand side of \eqref{eq:11} as
\[
\mathcal{L}(x_{k+1},y^{*})-\mathcal{L}(x^{*},y_{k+1}).
\tag{13}
\]

\item[Step 6.] Substitute \eqref{eq:12}–\eqref{eq:13} into \eqref{eq:11} and rearrange:
\[
\begin{aligned}
\Phi_{k}-\Phi_{k+1}
&\ge \mathcal{L}(x_{k+1},y^{*})-\mathcal{L}(x^{*},y_{k+1})
   +\underbrace{\frac{1}{2\tau}\|x_{k+1}-x_{k}\|^{2}
               +\frac{1}{2\sigma}\|y_{k+1}-y_{k}\|^{2}}_{\ge 0}\\
&\quad+\langle K(x_{k+1}-\bar{x}_{k}),y_{k+1}\rangle .
\end{aligned}
\tag{14}
\]

\item[Step 7.] By the definition of $\bar{x}_{k}=x_{k}+\theta(x_{k}-x_{k-1})$,
\[
x_{k+1}-\bar{x}_{k}=x_{k+1}-x_{k}-\theta(x_{k}-x_{k-1})
\]
and therefore, using Cauchy–Schwarz and the step‑size condition (A1),
\[
\bigl|\langle K(x_{k+1}-\bar{x}_{k}),y_{k+1}\rangle\bigr|
\le \frac{1}{2\tau}\|x_{k+1}-\bar{x}_{k}\|^{2}
    +\frac{\tau}{2}\|K^{\!\top}y_{k+1}\|^{2}
\le \frac{1}{2\tau}\|x_{k+1}-\bar{x}_{k}\|^{2}
    +\frac{1}{2\sigma}\|y_{k+1}-y_{k}\|^{2},
\tag{15}
\]
where the last inequality follows from $\tau\sigma\|K\|^{2}<1$.

\item[Step 8.] Plug \eqref{eq:15} into \eqref{eq:14} and cancel the non‑negative terms
$\frac{1}{2\sigma}\|y_{k+1}-y_{k}\|^{2}$. We obtain the fundamental descent inequality
\[
\Phi_{k}-\Phi_{k+1}
\ge \mathcal{L}(x_{k+1},y^{*})-\mathcal{L}(x^{*},y_{k+1})
    -\frac{1}{2\tau}\|x_{k+1}-\bar{x}_{k}\|^{2}.
\tag{16}
\]

\item[Step 9.] Using the definition of $\bar{x}_{k}=x_{k}+\theta(x_{k}-x_{k-1})$ and the bound
$\|x_{k+1}-\bar{x}_{k}\|^{2}\le 2\|x_{k+1}-x_{k}\|^{2}+2\theta^{2}\|x_{k}-x_{k-1}\|^{2}$,
the last term can be absorbed into the non‑negative quadratic terms that already appear in $\Phi_{k}-\Phi_{k+1}$. Hence we can drop it, yielding the simpler bound
\[
\Phi_{k}-\Phi_{k+1}
\ge \mathcal{L}(x_{k+1},y^{*})-\mathcal{L}(x^{*},y_{k+1}).
\tag{17}
\]

\item[Step 10.] Sum \eqref{eq:17} for $k=0,\dots,T-1$:
\[
\sum_{k=0}^{T-1}\bigl[\mathcal{L}(x_{k+1},y^{*})-\mathcal{L}(x^{*},y_{k+1})\bigr]
\le \Phi_{0}-\Phi_{T}
\le \Phi_{0},
\tag{18}
\]
because $\Phi_{T}\ge 0$.

\item[Step 11.] By convexity of $f$ and $g^{*}$, the saddle function $\mathcal{L}(\cdot,\cdot)$ is convex in its first argument and concave in its second. Hence
\[
\frac{1}{T}\sum_{k=1}^{T}\mathcal{L}(x_{k},y^{*})
   \ge \mathcal{L}\bigl(\bar{x}_{T},y^{*}\bigr),\qquad
\frac{1}{T}\sum_{k=1}^{T}\mathcal{L}(x^{*},y_{k})
   \le \mathcal{L}\bigl(x^{*},\bar{y}_{T}\bigr).
\tag{19}
\]

\item[Step 12.] Divide \eqref{eq:18} by $T$ and combine with \eqref{eq:19}:
\[
\mathcal{L}(\bar{x}_{T},y^{*})-\mathcal{L}(x^{*},\bar{y}_{T})
\le \frac{\Phi_{0}}{T}
   =\frac{1}{2\tau T}\norm{x_{0}-x^{*}}^{2}
    +\frac{1}{2\sigma T}\norm{y_{0}-y^{*}}^{2}.
\tag{20}
\]

\item[Step 13.] The left‑hand side of \eqref{eq:20} is precisely the primal–dual gap
$\mathcal{G}_{T}$ defined in \eqref{eq:2}. Therefore,
\[
\mathcal{G}_{T}\le\frac{C}{T},
\qquad
C:=\frac{1}{2\tau}\norm{x_{0}-x^{*}}^{2}
   +\frac{1}{2\sigma}\norm{y_{0}-y^{*}}^{2},
\]
which completes the proof. \qed
\end{enumerate}

% -------------------------------------------------
\section*{Rate Derivation (With Constants)}
From \eqref{eq:20} we see that the constant $C$ depends only on the initial distance to the saddle point and the chosen step sizes. Choosing the maximal admissible product $\tau\sigma=\frac{1}{\|K\|^{2}}(1-\varepsilon)$ with a small $\varepsilon>0$ balances the two terms in $C$ and yields the smallest possible bound.

% -------------------------------------------------
\section*{Special Cases}
\begin{enumerate}
  \item \textbf{Smooth $f$, $g\equiv0$:} $K=I$, $g^{*}=0$. The algorithm reduces to the classical gradient descent with step $\tau$, and the bound \eqref{eq:2} recovers the standard $O(1/T)$ rate.
  \item \textbf{Strongly convex $f$:} If $f$ is $\mu$‑strongly convex, one can strengthen \eqref{eq:9} to obtain a linear convergence factor $(1-\mu\tau)$ in place of the $O(1/T)$ term.
  \item \textbf{ADMM limit:} Setting $\theta=1$ and $\sigma=\tau$ with $K$ invertible, PDHG becomes equivalent to the classical Alternating Direction Method of Multipliers; the same $O(1/T)$ ergodic rate follows.
\end{enumerate}

\end{document}